\documentclass[11pt]{article}

\usepackage[margin=1in]{geometry}
\usepackage{fancyhdr}
\usepackage{enumitem}
\usepackage{listings}
\usepackage{verbatim}
\usepackage{setspace}

\newcommand{\code}[1]{\texttt{\detokenize{#1}}}

\pagestyle{fancy}
\fancyhf{}
\rhead{CPSC 250}
\lhead{Test 2}
\rfoot{\thepage}

\begin{document}

\begin{center}
{\Large \textbf{CPSC 250 -- Python for Data Manipulation}}\\
\vspace{0.2cm}
{\Large \textbf{Test 2}}\\
\vspace{0.5cm}
Name:\rule{3in}{0.4pt}
\end{center}

\vspace{0.25cm}

\textbf{Topics Covered:}
Functions, Parameter Passing, Classes and Objects, Constructors,
Collections of Objects, and Operator Overloading

\vspace{0.25cm}

\textbf{Total Points: 70}

\hrule

\vspace{0.3cm}

\section*{Part I -- Multiple Choice (10 points)}

Circle the best answer. Each question is worth 2 points.

\begin{enumerate}

\item What does the keyword \code{self} refer to inside a class?

\begin{enumerate}[label=\Alph*.]
\item The class itself
\item The current object
\item A constructor
\item A list of objects
\end{enumerate}

\vspace{0.2cm}

\item What is the purpose of the \code{__init__()} method?

\begin{enumerate}[label=\Alph*.]
\item Print an object
\item Compare two objects
\item Initialize an object when it is created
\item Delete an object
\end{enumerate}

\vspace{0.2cm}

\item What is the purpose of the \code{__str__()} method?

\begin{enumerate}[label=\Alph*.]
\item Create an object
\item Return a string representation of an object
\item Compare two objects
\item Modify an object
\end{enumerate}

\vspace{0.2cm}

\item If a class defines an \code{__eq__()} method, which operator uses it?

\begin{enumerate}[label=\Alph*.]
\item \code{+}
\item \code{*}
\item \code{==}
\item \code{=}
\end{enumerate}

\end{enumerate}

\newpage

\begin{enumerate}
\setcounter{enumi}{4}

\item Which statement creates a new object named \code{book}?

\begin{enumerate}[label=\Alph*.]
\item \code{book = Book}
\item \code{Book = book()}
\item \code{book = Book()}
\item \code{create Book()}
\end{enumerate}

\end{enumerate}

\newpage

\section*{Part II -- Find the Error (15 points)}

Each code segment contains one or more errors. Identify and correct them.

\vspace{0.3cm}

\textbf{6.}

\begin{verbatim}
class Student:
    def __init__(name):
        self.name = name
\end{verbatim}

\vspace{1.5in}

\textbf{7.}

\begin{verbatim}
class Student:
    def __init__(self, name):
        self.name = name

s = Student
print(s.name)
\end{verbatim}

\vspace{1.5in}

\textbf{8.}

\begin{verbatim}
students = []

students.append(Student)

print(students[0].name)
\end{verbatim}

\vspace{1.5in}

\newpage

\section*{Part III -- Code Writing (20 points)}

\textbf{9.} Write a class called \code{Book} with the following:

\begin{itemize}
\item \code{title}
\item \code{author}
\item constructor, \code{__init__()}
\item \code{__str__()} method
\end{itemize}

\vspace{5in}

\newpage

\textbf{10.} Given the \code{Book} class from Question 9, write code that:

\begin{itemize}
\item Creates three \code{Book} objects
\item Stores them in a list called \code{library}
\item Uses a loop to print all books
\end{itemize}

\vspace{6in}

\newpage

\section*{Part IV -- Code Commentary (15 points)}

Write comments describing what each section of code does.

\vspace{0.3cm}

\begin{verbatim}
# _______________________________________

class Student:

# _______________________________________

    def __init__(self, name):
        self.name = name

# _______________________________________

students = []

# _______________________________________

students.append(Student("Alice"))

# _______________________________________

for s in students:
    print(s)
\end{verbatim}

\vspace{2in}

\newpage

\section*{Part V -- Short Answer (10 points)}

\textbf{11.} Explain the difference between a class and an object.

\vspace{2.5in}

\textbf{12.} Explain why a list of objects is often more useful than creating many separate variables such as

\begin{verbatim}
student1
student2
student3
\end{verbatim}

Give at least one advantage.

\vspace{2.5in}

\hrule

\vspace{0.2cm}

\begin{center}
\textbf{End of Test}
\end{center}

\end{document}
